% !TeX root = main.tex
\chapter{Supplementary Design Material}\label{app:design}
\section{Terminal Feedback and Operator Controls}

The following rules refine the proposed scanner feedback schedule and extend it
to both attendance configurations and maintenance operations. The schedule
forms the implemented software contract. Host simulations and firmware builds
validate the control logic; physical mode switching and assembled LED/buzzer
timings still require bench confirmation. These checks do not constitute an
end-to-end website attendance validation. All indications
are transient. Green or red indications last 700~ms; a short operation beep is
100~ms, a key click is 20~ms and an enrolment-completion beep is 700~ms. Error and
retry signals contain two 120~ms beeps separated by 120~ms. Deliberate deletion
contains two 100~ms beeps separated by 100~ms. The scheduler is nonblocking and
key clicks do not interrupt an active success/error sequence.

% !TeX root = ../main.tex
\begingroup
\singlespacing\small
\setlength{\tabcolsep}{3pt}
\begin{longtable}{@{}>{\RaggedRight\arraybackslash}p{\dimexpr.24\linewidth-2\tabcolsep\relax}>{\RaggedRight\arraybackslash}p{\dimexpr.34\linewidth-2\tabcolsep\relax}>{\RaggedRight\arraybackslash}p{\dimexpr.105\linewidth-2\tabcolsep\relax}>{\RaggedRight\arraybackslash}p{\dimexpr.105\linewidth-2\tabcolsep\relax}>{\RaggedRight\arraybackslash}p{\dimexpr.21\linewidth-2\tabcolsep\relax}@{}}
\caption{Defined feedback for terminal events}\label{tab:terminal-feedback}\\
\toprule
\textbf{Event} & \textbf{Screen or result} & \textbf{Green} & \textbf{Red} & \textbf{Buzzer} \\
\midrule\endfirsthead
\toprule
\textbf{Event} & \textbf{Screen or result} & \textbf{Green} & \textbf{Red} & \textbf{Buzzer} \\
\midrule\endhead
\midrule\multicolumn{5}{r}{Continued on next page}\\\endfoot
\bottomrule\endlastfoot
Startup & Starting, then Home; attendance stopped & Off & Off & One short beep \\ \addlinespace
Enrolment attempt 1 accepted & Lift Finger, then Enrol 2 of 3 & Off & Off & None \\ \addlinespace
Enrolment attempt 2 accepted & Lift Finger, then Enrol 3 of 3 & Off & Off & None \\ \addlinespace
Enrolment complete & Enrolled; verified scanner slot & On & Off & One long beep \\ \addlinespace
Unreadable enrolment scan & Try Again; lift finger first & Off & On & Two error beeps \\ \addlinespace
Enrolment mismatch, full gallery or unconfirmed save & Specific failure; no success claim & Off & On & Two error beeps \\ \addlinespace
Already enrolled & Already Enrolled; existing slot & Off & On & Two error beeps \\ \addlinespace
Immediate attendance saved & Attendance Saved; local record saved & On & Off & One short beep \\ \addlinespace
Already marked & Already Checked In; same template/date & Off & On & Two error beeps \\ \addlinespace
Not recognised or unreadable attendance & No Match or Scan Failed; lift and retry & Off & On & Two error beeps \\ \addlinespace
Match without valid clock/storage & Matched / Not Saved; reason displayed & Off & On & Two error beeps \\ \addlinespace
Deferred capture saved & Capture Saved / Pending Match & Off & Off & One short beep \\ \addlinespace
Deferred capture unreadable & Try Again; lift and reposition & Off & On & Two error beeps \\ \addlinespace
Deferred communication, storage, integrity or capacity fault & Capture Stopped; files retained & Off & On & Two error beeps \\ \addlinespace
Held finger & Lift between scans; no extra record & Off & Off & None \\ \addlinespace
Missing class batch & Capture Not Started; prepare class via USB & Off & On & Two error beeps \\ \addlinespace
Mode confirmation & Use Immediate/Deferred? Existing records kept & Off & Off & Key click only \\ \addlinespace
Menu held after USB maintenance & Mode selection; capture remains stopped & Off & Off & One short beep \\ \addlinespace
Mode confirmed/start succeeds & Selected mode running & Off & Off & One short beep \\ \addlinespace
Physical stop & Selected mode stopped & Off & Off & One short beep \\ \addlinespace
Batch prepared/selected & Batch prepared; class UUID available & Off & Off & One short beep \\ \addlinespace
Menu/cancel & Return without saving or deleting & Off & Off & Click; cancel gives one short beep \\ \addlinespace
Confirmation or action timeout & Cancelled, or Action Timed Out & Off & Off/On & Cancel beep, or two error beeps \\ \addlinespace
Delete awaiting confirmation & Identify, lift, then fresh Delete press & Off & Off & Key click only \\ \addlinespace
Delete or full-library erase verified & Deleted / Library Erased; history kept & On & Off & Two deletion beeps \\ \addlinespace
Delete/erase unconfirmed & Unconfirmed; inspect library & Off & On & Two error beeps \\ \addlinespace
Clock saved & Clock Saved; readback verified & Off & Off & One short beep \\ \addlinespace
Clock save failed & Clock Not Saved; check DS1307 & Off & On & Two error beeps \\ \addlinespace
Log export complete & Log Exported; USB output complete & Off & Off & One short beep \\ \addlinespace
Deferred export complete & Batch Exported; device files retained & Off & Off & One short beep \\ \addlinespace
Export failure & Specific failure; evidence retained & Off & On & Two error beeps \\ \addlinespace
Diagnostics, backup or restore & Attendance Paused; console protocol results & Off & Off & Silent during transfers \\ \addlinespace
Confirmed USB filesystem format & Console completion/failure; all local files affected & Off & Off & Silent maintenance operation \\ \addlinespace
Background missing clock, sensor or storage & Home readiness warning & Off & Off & No repeated alarm; attempted operation reports failure \\ \addlinespace
Button press & Current screen retained unless action changes it & Off & Off & Short click; active outcome feedback has priority \\
\end{longtable}
\endgroup


Green indicates a verified local persistent biometric operation, not a generic
acknowledgement and not a server-side identity guarantee. Deferred captures
remain unidentified even when stored correctly. A duplicate immediate check-in
does not add another row. Export success confirms the local stream operation,
not a durable laptop file or backend receipt. USB protocol maintenance suppresses
automatic capture and uses console-specific results without incidental green
indications or unsolicited feedback inside transfer frames.

\section{Mode Selection and Recovery}

The installed application starts with attendance stopped. On Attendance Mode,
Enrol selects Immediate and Delete selects Deferred; a fresh Enrol press confirms
the selection and starts the chosen path when its prerequisites are satisfied.
Menu cancels. A confirmation expires after 20~s. No mode selection erases data.
The Start/Stop page changes acquisition state without changing attendance type.

Deferred Class Batch cycles through valid batches prepared by the USB tool and
shows the selected class-session UUID.\@ The terminal retains the selection but
requires explicit re-arming after restart. An interrupted or damaged batch is
not silently repaired, replaced or reformatted. An unavailable or full capture
store stops capture and retains existing evidence. Existing scanner diagnostic,
backup and restore commands pause acquisition until an explicit restart.

After USB maintenance, a fresh two-second Menu hold opens attendance-mode
selection with capture still stopped. The UI handles this hold only after the
active serial transaction has completed, preserving framed transfers. The
deferred statistics page counts pending captures rather than identified students;
its export page directs the operator to the laptop batch downloader. Unassigned
Delete presses on Start/Stop and Class Batch pages do not initiate deletion.

\section{States Reserved for Incomplete Integration}

Automatic student-enrolment upload and database binding remain unfinished. The
future laptop/website interface must distinguish Sending Enrolment, Registration
Pending and Registration Saved. Only a durable database acknowledgement may
produce the last state; sensor-only enrolment is insufficient. Likewise, deferred
matching requires separate Matched, Unknown, Ambiguous and Out-of-roster review
states before a student/class attendance write. These integrations remain
unfinished.

\section{Remaining Assembly Material}

The final circuit, resistor values and measured driver-stage behaviour still
require assembly evidence. The existing conceptual and methodology diagrams are
included in Chapters Two and Three; their editable sources remain in folder/.
